\chapter{完整源码清单：算法题解验证器}

本附录把 \filepath{books/algorithm-interview/labs/cpp20-examples} 的主线工程完整放入书中。阅读顺序是：先看 CMake 如何组织库和测试，再看头文件公开了哪些题解函数，然后读实现文件，最后读测试文件确认每个优化都和 reference 对照。这里的清单来自真实工程文件，不是手抄节选。

\section{一、CMakeLists.txt}

这个文件定义算法实践库、测试目标和依赖关系。阅读时重点看 include 目录如何暴露给目标、实现文件如何进入库、GTest 如何接入测试；如果以后新增题型，应该先让它进入这个构建图。

\lstinputlisting[language=bash,caption={cpp20-examples/CMakeLists.txt}]{../../labs/cpp20-examples/CMakeLists.txt}

\section{二、公共接口}

\filepath{array_patterns.hpp} 是数组、双指针、滑动窗口、前缀和与 Kadane 的公共接口。它只声明题解函数，不藏实现细节；读者应先从这里看每道题的输入输出合同。

\lstinputlisting[language=C++,caption={include/algobook/array\_patterns.hpp}]{../../labs/cpp20-examples/include/algobook/array_patterns.hpp}

\topic{图搜索公共头}

\filepath{graph_patterns.hpp} 只暴露 BFS 图搜索案例。重点是参数类型：网格按值传入，允许实现修改访问状态，不把副作用泄露给调用者。

\lstinputlisting[language=C++,caption={include/algobook/graph\_patterns.hpp}]{../../labs/cpp20-examples/include/algobook/graph_patterns.hpp}

\topic{回溯公共头}

\filepath{backtracking_patterns.hpp} 暴露 24 点判断函数。当前实现使用显式栈模拟搜索状态，避免运行时递归；读者应把它和正文里“选择、尝试、恢复”的解释对起来看。

\lstinputlisting[language=C++,caption={include/algobook/backtracking\_patterns.hpp}]{../../labs/cpp20-examples/include/algobook/backtracking_patterns.hpp}

\section{三、完整实现}

这个实现文件是算法书最重要的代码样本。每组函数通常先给 reference，再给优化版：两数之和从双循环到哈希表，容器盛水从枚举到双指针，子数组问题从枚举到窗口或前缀哈希，最大子数组从枚举到 Kadane。读代码时不要只看最终公式，要看每个函数保留了哪些边界检查。

\lstinputlisting[language=C++,caption={src/array\_patterns.cpp}]{../../labs/cpp20-examples/src/array_patterns.cpp}

\topic{图搜索完整实现}

图搜索实现集中训练 BFS 的两个核心边界：什么时候入队，什么时候标记访问。腐烂橘子用层数表达分钟，岛屿数量用一次 BFS 吃掉一个连通块。读者应检查不可达、新鲜数量、对角线不连通这些边界如何进入代码。

\lstinputlisting[language=C++,caption={src/graph\_patterns.cpp}]{../../labs/cpp20-examples/src/graph_patterns.cpp}

\topic{回溯完整实现}

24 点实现展示了如何把递归搜索改写成显式状态栈。每个状态保存当前剩余数字；每一步选两个数生成下一状态。读者重点看减法、除法的两个方向，以及接近 0 的除数如何被跳过。

\lstinputlisting[language=C++,caption={src/backtracking\_patterns.cpp}]{../../labs/cpp20-examples/src/backtracking_patterns.cpp}

\section{四、完整测试}

测试文件是题解进入书稿的验收门槛。它不仅跑样例，还保留重复值、负数、全 0、不可达、除法组合等边界。新增案例时，必须把 reference 与优化版放在同一批输入上对照。

\lstinputlisting[language=C++,caption={tests/algorithm\_book\_tests.cpp}]{../../labs/cpp20-examples/tests/algorithm_book_tests.cpp}
